\section{Conclusão}

A análise do paradigma da Computação Autônoma revela que a indústria de Tecnologia da Informação atingiu um ponto de inflexão onde a escalabilidade dos sistemas não pode mais ser dissociada da automação de sua gestão. A evolução histórica, que partiu de sistemas isolados para ecossistemas hiper-conectados em nuvem e dispositivos de Internet das Coisas (IoT), evidenciou que a capacidade cognitiva humana é o principal limitador para a operação segura e eficiente das infraestruturas contemporâneas. A ``crise da complexidade'', discutida desde o manifesto original da IBM \cite{ibm2001}, consolidou-se como o desafio central para engenheiros e pesquisadores da área.

Ao longo deste trabalho, observou-se que a implementação das propriedades \textit{Self-CH}:  auto-configuração, auto-cura, auto-otimização e auto-proteção, são fundamentadas na maturidade do ciclo de controle MAPE-K. Este modelo permite que o conhecimento do domínio seja codificado em gerentes autônomos capazes de agir proativamente frente a falhas ou flutuações de demanda. As aplicações discutidas, tanto no contexto de \textit{Cloud Computing} quanto em sistemas pervasivos, demonstraram que a autonomia não é apenas uma questão de conveniência operacional, mas de viabilidade econômica. A capacidade de reduzir o custo total de propriedade (TCO) e mitigar o erro humano é o que permite a sustentabilidade de serviços globais que operam sob rigorosos Acordos de Nível de Serviço (SLAs) \cite{sterritt2005, correa2007}.

Entretanto, o caminho para a plena autonomia ainda enfrenta barreiras significativas. A transição de sistemas que apenas auxiliam a decisão humana para sistemas que agem de forma totalmente independente exige avanços em termos de interoperabilidade e, crucialmente, de confiança. A estabilidade de algoritmos de controle em ambientes altamente dinâmicos e a garantia de que políticas de alto nível não resultem em estados sistêmicos indesejados permanecem como fronteiras de pesquisa ativa. Nesse sentido, a convergência da Computação Autônoma com técnicas de verificação formal e métodos rigorosos de modelagem apresenta-se como uma tendência promissora para assegurar que a auto-gestão seja tão previsível quanto eficiente \cite{lalanda2013, kephart2003}.

Em conclusão, a Computação Autônoma representa uma mudança fundamental na filosofia de design de software: passamos de sistemas que ``fazem o que mandamos'' para sistemas que ``fazem o que queremos''. À medida que avançamos para uma era de trilhões de dispositivos interconectados, a inteligência intrínseca proposta por este paradigma deixará de ser um diferencial competitivo para se tornar um requisito mandatório. O sucesso dessa jornada garantirá que a tecnologia continue a evoluir como uma extensão invisível e resiliente da atividade humana, capaz de se regenerar, se proteger e se otimizar de forma transparente e segura.